\chapter{Один комплексный базис и несколько ветвей сети}
\label{adv:chapter}

Допустим, для семидесяти значений параметра уже вычислены комплексные поля.
Можно ли по ним восстановить поле при новом значении параметра?
Для ответа проведём два опыта. Сначала дадим алгоритму само проверочное
поле и попросим подобрать его координаты в обучающем базисе. Затем оставим
только параметры задачи и поручим подбор координат сети. Разница между
результатами покажет, на каком этапе теряется точность.

\section{Какие поля получает алгоритм}
\label{adv:family}

Один снимок $u_q\in\C^N$ содержит $N$ коэффициентов рассеянного поля
в конечноэлементном базисе одной сетки. Номер $q$ указывает на пару
$(k_q,\theta_q)$, где $k$ --- волновое число, а $\theta$ --- направление
падающей волны. Геометрия и материалы одинаковы для всех снимков.
В наборе \texttt{mscale\_sno\_100} зафиксирован также угол
$\theta=0$. Следовательно, изменяется один параметр, и приближаемая
зависимость имеет вид $k\mapsto u_h(k,0)$.

Всего в серии 100 решений при $k\in[0{,}1;4]$. Из них 70 отведены для
обучения, 15 для валидации (выбора весов сохранённой эпохи) и 15 для
итогового теста. Каждая сравниваемая модель получает одни и те же задачи.
Среднее поле, базис, его ранг и параметры нормировки определяются по 70
обучающим снимкам. Крайние значения $k$ входят в обучение, поэтому тест
проверяет интерполяцию по $k$ в одной среде.

\section{Как получить комплексные моды из снимков}
\label{adv:pod}

Пусть $v_h$ и $u_h$ --- конечноэлементные поля с векторами коэффициентов
$v$ и $u$. Их скалярное произведение $\int_\Omega\overline{v_h}u_h\,dx$
равно $v^*Mu$. Здесь $M_{ij}=\int_\Omega\psi_i\psi_j\,dx$ ---
согласованная матрица масс вещественного базиса $\{\psi_i\}$,
а звёздочка обозначает сопряжённое транспонирование. Матрица $M$
эрмитова и положительно определена. Она учитывает площадь элементов
и перекрытие базисных функций. Поэтому $\norm{u}_M^2=u^*Mu$
измеряет квадрат $L^2$-нормы пространственного поля, в том числе на
неравномерной сетке.

Вычтем из каждого обучающего поля комплексное среднее $\mu$ и запишем
отклонения в столбцы матрицы $S$:
\[
 \mu=\frac1n\sum_{q=1}^n u_q,\qquad
 S=[u_1-\mu,\ldots,u_n-\mu]\in\C^{N\times n},\qquad n=70.
\]
Сумма столбцов равна нулю, поэтому их ранг не превосходит $69$.
Среднее сохраняется отдельно. Для поиска мод достаточно решить
собственную задачу размера $70\times70$:
\begin{equation}
 G=S^*MS=V\Lambda V^*,\qquad
 \Phi_R=SV_R\Lambda_R^{-1/2}.
 \label{adv:gram}
\end{equation}
Здесь $V_R$ содержит собственные векторы, соответствующие $R$ наибольшим
численно положительным собственным значениям $\lambda_j$.
Формула имеет простую проверку по размерам. Матрица $V_R$ размера
$n\times R$ задаёт смеси снимков, а $S V_R$ размера $N\times R$ уже
задаёт пространственные поля. Деление каждого столбца на
$\sqrt{\lambda_j}$ нормирует его. Подставив \eqref{adv:gram}, получим
$\Phi_R^*M\Phi_R=I_R$.

При точной арифметике $\lambda_j$ --- вклады мод в суммарную энергию
отклонений $u_q-\mu$. Код отбрасывает численные нули и ищет минимальное
$R$, для которого
\begin{equation}
 \frac{\sum_{j>R}\lambda_j}{\sum_j\lambda_j}\leq 10^{-8},
 \label{adv:energy}
\end{equation}
с верхним пределом $R\leq60$. Если предел достигнут раньше порога,
остаётся максимально допустимый ранг. В этом опыте хватило \textbf{17 мод}.
После построения базиса код корректирует его ортонормированность малым
разложением Холецкого. Условие \eqref{adv:energy} относится к сумме
квадратов отклонений обучающих полей. Относительную ошибку отдельного
нового поля ещё предстоит измерить.

\section{Проверка базиса до обучения сети}
\label{adv:loss}

Возьмём известное поле $u$ из тестовой выборки. Его наилучшие координаты
в метрике $M$, проекция и остаток равны
\[
 a_* = \Phi_R^*M(u-\mu),\qquad
 u_* = \mu+\Phi_R a_*,\qquad r=u-u_* .
\]
Проекция оставляет от $u-\mu$ ровно ту часть, которую можно выразить
через моды. Действительно, $\Phi_R^*Mr=0$. Для произвольного прогноза
$\widehat u=\mu+\Phi_R\widehat a$ из этого следует
\begin{equation}
 \norm{u-\widehat u}_M^2
 =\underbrace{\norm{r}_M^2}_{\text{остаток базиса}}
 +\underbrace{\norm{a_*-\widehat a}_2^2}_{\text{ошибка координат}}.
 \label{adv:pyth}
\end{equation}
Первое слагаемое одинаково для всех сетей с данным декодером. Второе
можно измерить, даже не восстанавливая $N$ значений прогноза.
При $u\ne0$ деление суммы на $\norm{u}_M^2$ и извлечение корня дают
относительную ошибку поля. При обучении усредняют эти ошибки по примерам.

\begin{worked}{Практикум. Проверьте проекцию с комплексным коэффициентом}
Пусть $M=I_2$, $\mu=0$, $\phi=(1,\ii)^T/\sqrt{2}$ --- единственная
сохранённая мода, а $u=(0,1)^T$. Найдите наилучшую реконструкцию и
ошибку сети, если сеть вернула $\widehat a=0$.

\textbf{Решение.} Сопряжение даёт $a_*=\phi^*u=-\ii/\sqrt{2}$, откуда
$u_*=(-\ii,1)^T/2$. Остаток $r=(\ii,1)^T/2$ имеет квадрат нормы $1/2$,
причём $\phi^*r=(\ii-\ii)/(2\sqrt{2})=0$. Ошибка координаты добавляет ещё $1/2$.
Поскольку $\norm{u}_2=1$, относительная ошибка наилучшей проекции равна
$1/\sqrt2$, а прогноза --- $1$. Поле $\widehat u=0$ даёт тот же ответ
при непосредственной проверке.
\end{worked}

Для проверки реализации удобно обработать сразу несколько проверочных
полей. В следующем фрагменте строка \texttt{U[q]} --- одно поле,
\texttt{Phi} имеет форму \texttt{(N, R)}, а \texttt{M} --- разреженная
матрица масс. В строке восстановления сопряжения нет.
\begin{lstlisting}
centered = U - mean
coeff = (M @ centered.T).T @ Phi.conj()
projected = mean + coeff @ Phi.T
residual = U - projected
gram = Phi.conj().T @ (M @ Phi)
orthogonality = (M @ residual.T).T @ Phi.conj()
\end{lstlisting}
Матрица \texttt{gram} должна быть близка к единичной, а
\texttt{orthogonality} --- к нулевой в пределах точности вычислений
и масштаба полей. Эти два теста помогают обнаружить перепутанные оси,
пропущенное сопряжение или матрицу масс от другой сетки.

\section{Как ветви предсказывают координаты}
\label{adv:architecture}

При новом запросе известны только $k$ и $\theta$. Сеть получает вектор
$x=((k-\bar k)/s_k,\cos\theta,\sin\theta)$, где $\bar k$ и $s_k$ ---
обучающие среднее и стандартное отклонение $k$. В рассматриваемой серии
последние два признака равны 1 и 0 для всех примеров.

Ветвь $j$ умножает все три признака на свой обучаемый скаляр $c_j$.
Её первый слой переводит $c_j x$ в $R W$ чисел. Мы укладываем их в
таблицу $Z$ с $R$ строками по числу мод и $W$ столбцами по числу скрытых
признаков. Уже после первого слоя у каждой моды своя строка. После
нелинейности применяются $D$ блоков вида
\[
 Z\ \longleftarrow\ Z+D^{-1/2}
 \rho\!\left(ZL+\mathbf1 b^T + BZC\right).
\]
Матрицы $L,C\in\R^{W\times W}$ преобразуют признаки, $b\in\R^W$ задаёт
смещение, а $B\in\R^{R\times R}$ смешивает строки разных мод. У каждого
блока и ветви свои веса. В многоветвевой модели $\rho=\sin$.
Ветви начинают обучение с разными $c_j$, поэтому изменение $k$ подаётся
на их входы в разных масштабах.
Пространственное разрешение во всех ветвях задаётся одной матрицей $\Phi_R$.

Выходной слой каждой ветви превращает строку из $W$ чисел в пару
$(\operatorname{Re}p_{jr},\operatorname{Im}p_{jr})$ для моды $r$.
Результат всей сети и поле имеют вид
\[
 s_a=\sqrt{\frac{1}{nR}\sum_{q=1}^n\sum_{r=1}^R|a_{qr}|^2},\qquad
 \widehat a_r=s_a\sum_{j=1}^J\gamma_j p_{jr},\qquad
 \widehat u=\mu+\Phi_R\widehat a.
\]
Единый масштаб $s_a$ возвращает нормированные выходы в физические единицы.
Веса $\gamma_j$ изначально равны $1/J$ и затем обучаются без ограничения
на сумму. Среднее $\mu$ добавляется один раз, после сложения ветвей.

\begin{worked}{Проверка размеров и числа параметров}
Подсчитаем веса одной ветви. Первый слой с тремя входами и смещением
требует $4RW$ параметров. Один блок содержит $2W^2+W+R^2$, выходной
слой --- $2W+2$. Добавим $c_j$ и $\gamma_j$. При $R=17$, $W=8$, $D=4$
получаем $544+4(128+8+289)+18+2=2264$ параметра.
У восьми ветвей их $18\,112$, а выход по-прежнему состоит из $17$
комплексных коэффициентов. Сеть с одной ветвью и подобранной шириной
в этом сравнении имеет $18\,380$ параметров.
\end{worked}

\clearpage
\section{Что показали пятнадцать тестовых полей}
\label{adv:comparison}

В таблице~\ref{adv:results} все три запуска используют один базис и те же
15 тестовых задач. Номер запуска служит начальным значением генератора,
от которого зависят начальные веса и порядок примеров при обучении.
В каждом запуске взяты веса эпохи с наименьшей средней ошибкой
на валидации. Числа в таблице --- проценты относительной комплексной
$L^2$-ошибки рассеянного поля по формуле~\eqref{exp:metric}. Для каждой строки
сначала вычислена ошибка каждого снимка, затем взяты соответствующие статистики.

\begin{table}[htbp]
\centering
\caption{Восьмиветвевая модель. Меньшее значение означает более точное поле.}
\label{adv:results}
\begin{tabular}{@{}lrrr@{}}
\toprule
Номер запуска & Среднее, \% & 95-й процентиль, \% & Максимум, \%\\
\midrule
11 & 65,85 & 84,55 & 85,26\\
29 & 60,96 & 87,84 & 89,87\\
47 & 60,19 & 94,58 & 96,14\\
\bottomrule
\end{tabular}
\end{table}

При координатах из истинных тестовых полей средняя ошибка того же
декодера равна всего $0{,}01579\%$. Значит, пространство ранга 17
хорошо представляет именно эти поля. Разрыв с таблицей указывает на
неточность предсказанных координат. Из чисел нельзя установить, как
разделить её между влиянием объёма выборки и процедурой обучения.

Для сравнения возьмём одну ветвь с той же синусоидальной активацией
и близким числом параметров. При начальных значениях 11, 29 и 47
её средние ошибки составили соответственно $71{,}15\%$, $74{,}99\%$
и $69{,}64\%$. Во всех трёх запусках восемь ветвей уменьшили среднюю
ошибку. Эти три наблюдения относятся к одной и той же тестовой выборке.

Общий эталон для этих запросов --- FEM-поля на сетке с $h\approx0{,}02$.
На той же квадратуре средняя ошибка FEM с шагом $h=0{,}04$ составила
$1{,}0170\%$, а у четырёх вариантов GMsFEM она лежала между $6{,}8244\%$
и $31{,}3606\%$. Средняя ошибка каждого запуска сети оказалась выше.
Все эти сравнения относятся к одним и тем же 15 физическим запросам.

\begin{worked}{Вопрос по таблице}
Верно ли, что запуск 47 дал самый точный результат для каждой частоты?
\textbf{Решение.} У запуска 47 наименьшее среднее, но наибольший
максимум. На частоте, где его ошибка достигла $96{,}14\%$, оба других
запуска дали меньшую ошибку: даже их максимумы ниже. Чтобы узнать,
на каких ещё частотах меняется порядок, нужны попарные ошибки по снимкам.
\end{worked}

Данные для воспроизведения таблицы находятся в
\path{experiments/mscale_sno_100/results/summary.csv}, а проекционные
остатки --- в \path{results/training_summary.json} той же серии.
Модули \texttt{pod.py} и \texttt{model.py} с реализациями базиса и ветвей
находятся в \path{experiments/gms_sno_100/}.

\clearpage
\section{Что проверяет смена сетки}
\label{adv:mesh}

\textbf{Вычисление уже полученного поля в новых точках.}
Для непрерывного кусочно-линейного поля (P1) пусть $T$ переводит
исходные узловые значения в значения в новых точках внутри сетки
с помощью барицентрической интерполяции. Линейность даёт
\[
 T(\mu+\Phi_R a)=T\mu+(T\Phi_R)a .
\]
Значит, среднее и моды можно перенести один раз, а затем подставлять
новые коэффициенты. Равенство правой и левой частей --- проверка порядка
узлов и корректности реализации переноса. Оно выполняется даже при
большой ошибке исходного поля относительно решения уравнения.

\textbf{Сравнение решений на двух сетках.}
Здесь нужны два решения одной физической задачи. Для исследования влияния
сетки на весь метод заново строят FEM-метки, базис и модель на каждой
сетке при одинаковых физических запросах. Следует сравнивать декодированные
поля в одних физических точках и интегрировать разность на общей
проверенной квадратуре. Даже на одной сетке моду можно домножить на
$e^{\ii\alpha}$ и компенсировать фазу в коэффициенте. Поэтому разность
модальных координат из двух независимых базисов сама по себе малоинформативна.

В сеточном пилоте проекта для $k=20$ различия квадратичных FEM-решений
между последовательными уровнями составили $5{,}40\%$ и $0{,}50\%$.
В этих числах участвуют только FEM-поля. Для вывода о сеточной
устойчивости сети нужны предсказания на обеих сетках и их сравнение
с общим эталоном. Особенно важно разделять эти проверки при повышении
$k$, когда пространственная сетка должна разрешать более короткие волны.

\textbf{Изменение диапазона параметра.}
Для новой обучающей серии могут измениться и формы полей, и их нормы.
Например, в высокочастотной серии проекта ранг выбирается по
среднеквадратической относительной ошибке проекции на обучающих полях,
$\big(n^{-1}\sum_q \|u_q-u_{*,q}\|_M^2/\|u_q\|_M^2\big)^{1/2}$.
Здесь каждое поле входит со своим знаменателем. Доля общей
центрированной энергии в \eqref{adv:energy} сильнее зависит от снимков
с большими отклонениями от среднего. Правило выбора ранга и метрику
проверки следует указывать вместе с числом оставленных мод.
